Fix $d\ge 1$.
For a finite set $A\subseteq\mathbb{R}^d$ of size $n$, define its diameter
\[\operatorname{diam}(A):=\max\{\|x-y\|: x,y\in A\}.
\]
Assume that for any two distinct unordered pairs $\{x,y\}\ne\{x',y'\}$ with $x\ne y$ and $x'\ne y'$, the corresponding distances satisfy
\[\big|\|x-y\| - \|x'-y'\|\big|\ge 1.
\]
(Equivalently: the multiset of $\binom{n}{2}$ pairwise distances is $1$-separated.)
The conjecture asks whether necessarily
\[
\operatorname{diam}(A) \ge (1+o(1))\, n^2\quad\text{as } n\to\infty.
\]

\medskip
\noindent\emph{Potential ambiguity.}
The phrase ``all pairwise distances differ by at least $1$'' is interpreted above as a separation condition between distances of \emph{distinct pairs}. This is consistent with the statement that the trivial diameter lower bound is $\binom n2$ (see Lemma 670.2).
